\section{Introduction}\label{sec:intro}
Network intrusion detection systems (NIDS) look at network flows and decide whether each one is benign or one of
several attack types. A flow is a single row of summary statistics for one connection. Recent work feeds such rows to
language models by writing them out as text~\cite{hegselmann2022tabllm,borisov2022great}.
SecurityBERT~\cite{ferrag2024securitybert} goes one step further. Its Privacy-Preserving Fixed-Length Encoding (PPFLE)
replaces each \texttt{column\$value} pair with a truncated hash, and a small BERT~\cite{devlin2019bert} is trained on
the resulting tokens. The paper reports \res{exp1.paper.accuracy} accuracy on the \res{nclasses.edge_iiotset}-class
Edge-IIoTset~\cite{ferrag2022edge} and claims that the hashes protect the raw values.

Three questions remain open. First, the encoding itself is never tested. Hashing removes both the numeric order and
any English meaning, and it is unclear whether a readable encoding would do better. Second, the baselines were weak.
On tabular data, a tuned gradient-boosted tree is the natural competitor. Third, real deployments meet networks they
were not trained on, and the privacy claim has never been checked. We therefore ask:
\begin{description}[leftmargin=3.2em, labelwidth=2.6em, font=\bfseries]
  \item[RQ1] Does the PPFLE + small BERT result reproduce on Edge-IIoTset?
  \item[RQ2] Within one dataset, how do four encodings compare with each other and with XGBoost?
  \item[RQ3] When a model is trained on one dataset and tested on another, which encoding generalises best?
  \item[RQ4] Is PPFLE actually privacy-preserving?
\end{description}
We also test four hypotheses:
\begin{description}[leftmargin=3.2em, labelwidth=2.6em, font=\bfseries, itemsep=0pt]
  \item[H1] PPFLE does well within a dataset.
  \item[H2] XGBoost is as good or better within a dataset.
  \item[H3] Readable or binned encodings transfer better than raw PPFLE.
  \item[H4] A pretrained model benefits from readable encodings, but not from hashes.
\end{description}

We make four contributions. We run a controlled benchmark of encodings and language models against XGBoost
(\cref{sec:results}). We find a label leak in Edge-IIoTset that explains the published accuracy (\cref{sec:leak}). We
show that a simple dictionary attack reverses most PPFLE tokens (\cref{sec:privacy}). Finally, we discuss when the NLP
route pays off (\cref{sec:discussion}). The box below gives the short answers.

\begin{findings}[Key findings]
\begin{description}[leftmargin=3.2em, labelwidth=2.6em, font=\bfseries, itemsep=2pt]
  \item[RQ1] The published accuracy reproduces, but only because of a label leak. PPFLE + small BERT reaches
    \res{exp1.bert_small.e1_ppfle.edge_iiotset_raw.multiclass.accuracy} on the data as published (paper:
    \res{exp1.paper.accuracy}) and \res{exp1.bert_small.e1_ppfle.edge_iiotset.multiclass.accuracy} once the leak is
    removed.
  \item[RQ2] XGBoost is as good as or better than every text model on every dataset (H2: \res{hyp.H2.verdict}).
    Raw PPFLE is never the best encoding (H1: \res{hyp.H1.verdict}).
  \item[RQ3] Binning before hashing transfers best. For the small BERT, off-diagonal macro-F1 rises from
    \res{exp3.bert_small.e1_ppfle.off.macro_f1} to \res{exp3.bert_small.e2_ppfle_bins.off.macro_f1}, and DistilBERT
    with binned PPFLE is the best method overall at \res{exp3.distilbert.e2_ppfle_bins.off.macro_f1} (H3:
    \res{hyp.H3.verdict}). Within a dataset, pretraining helps on readable text. Across datasets it helps on
    hashes instead (H4: \res{hyp.H4.verdict}).
  \item[RQ4] PPFLE is not private. A dictionary attack reverses
    \res{privacy.e1_ppfle.public.token_share.min}--\res{privacy.e1_ppfle.public.token_share.max}\,\% of the NetFlow
    test tokens and \res{privacy_edge.edge_iiotset_raw.e1_ppfle.bruteforce.token_share}\,\% on Edge-IIoTset.
\end{description}
\end{findings}
% CHECK hyp.H1-H4 verdicts and the E2 ranking in exp3_cross
